% SOURCE_AUTHORITY: sga5_fr_workpass.tex, lines 631--679 (SHA256 791F4EFFC5E02832D5D77ED03518C8156D6F07E4C8238B03545DB93D883FBB28; LF-normalized unit).
\subsection*{Lema 2.2 (cf. [H] V 3.3)}
Sea $X$ un esquema localmente noetheriano y conexo, y sean $P, P' \in \ob \D_c^{-}(X)$ tales que $P \mathop{\otimes}\limits^{\mathrm{L}} P' \xrightarrow{\sim} A_X$. Existen entonces un $A_X$-módulo invertible $L$ y un entero $r$ tales que
\[
P \xrightarrow{\sim} L[r], \quad P' \xrightarrow{\sim} L^{\vee}[-r].
\]

La demostración se hará en dos pasos.

\textbf{1) Caso particular:} $X = \Spec(k)$, con $k$ un cuerpo separablemente cerrado. Basta reproducir la demostración (por lo demás sencilla) dada en [H] (loc.~cit.). El punto importante es que (EGA~$0_{\mathrm{I}}$~5.4.3), si $M$ y $M'$ son dos $A$-módulos de tipo finito tales que $M \otimes M' \xrightarrow{\sim} A$, entonces $M$ y $M'$ son invertibles y $M' \xrightarrow{\sim} M^{\vee}$; es aquí donde interviene la hipótesis de que $A$ sea local.

\textbf{2) Caso general.}

En virtud del caso particular precedente, para todo $x \in X$ existe un entero $r(x)$ tal que $P_{\ol{x}} \xrightarrow{\sim} A_{\ol{x}}[r(x)]$, $P'_{\ol{x}} \xrightarrow{\sim} A_{\ol{x}}[-r(x)]$, es decir, $\uH^i(P_{\ol{x}}) = 0$ (respectivamente, $\uH^i(P'_{\ol{x}}) = 0$) para $i \neq -r(x)$ (respectivamente, $+r(x)$), y $\uH^{-r(x)}(P_{\ol{x}}) \xrightarrow{\sim} A_{\ol{x}}$, $\uH^{+r(x)}(P'_{\ol{x}}) \xrightarrow{\sim} A_{\ol{x}}$. Se trata de demostrar que la función $x \mapsto r(x)$ es localmente constante. En efecto, supongamos demostrado este punto. Como $X$ es conexo, $r$ es constante, digamos de valor $r_0$. Se tiene entonces $\uH^i(P) = 0$ para $i \neq -r_0$, $\uH^i(P') = 0$ para $i \neq r_0$, y $\uH^{-r_0}(P) \otimes \uH^{r_0}(P') \xrightarrow{\sim} A$. Pongamos $\uH^{-r_0}(P) = L$, $\uH^{r_0}(P') = L'$.

Sea $u : \ol{x} \to \ol{y}$ una flecha de especialización (SGAA~VIII~7.2), y sean $u^{*} : L_{\ol{y}} \to L_{\ol{x}}$, $u'^{*} : L'_{\ol{y}} \to L'_{\ol{x}}$ los morfismos de especialización correspondientes (SGAA~VIII~7.7). Se tiene un diagrama conmutativo
\begin{equation}
\begin{tikzcd}
L_{\ol{y}} \otimes L'_{\ol{y}} \arrow[r,"\sim"] \arrow[d,"u^{*} \otimes u'^{*}"'] & A_{\ol{y}} \arrow[d] \\
L_{\ol{x}} \otimes L'_{\ol{x}} \arrow[r,"\sim"] & A_{\ol{x}}
\end{tikzcd}
\tag{*}
\end{equation}
Elijamos una base $e$ (respectivamente, $f$) de $L_{\ol{x}}$ (respectivamente, $L_{\ol{y}}$) sobre $A_{\ol{x}}$ (respectivamente, $A_{\ol{y}}$), y denotemos por $e'$ (respectivamente, $f'$) la base ``dual'' de $L'_{\ol{x}}$ (respectivamente, $L'_{\ol{y}}$), es decir, aquella para la cual la imagen de $e \otimes e'$ (respectivamente, $f \otimes f'$) en $A_{\ol{x}}$ (respectivamente, $A_{\ol{y}}$) es igual a 1. Se tiene $u^{*}(f) = \lambda e$, $u'^{*}(f') = \lambda' e'$, y la conmutatividad de $(*)$ implica $\lambda\lambda' = 1$; por tanto, $u^{*}$ es un isomorfismo (y $u'^{*}$ el isomorfismo ``contragrediente''). En virtud de (SGAA~IX~2.13), esto implica que $L$ y $L'$ son localmente constantes y, por tanto, de presentación finita como $A_X$-módulos. Por consiguiente, si $x \in X$, la flecha canónica $\underline{\Hom}(L,F)_{\ol{x}} \to \Hom(L_{\ol{x}}, F_{\ol{x}})$ es un isomorfismo para todo $A_X$-módulo $F$. Como las fibras geométricas de $L$ son invertibles, se sigue que $L$ es invertible. Del mismo modo, $L'$ es invertible y $L' \xrightarrow{\sim} L^{\vee}$. Por tanto, se tiene $P \xrightarrow{\sim} L[-r_0]$, $P' \xrightarrow{\sim} L^{\vee}[+r_0]$.

Todo se reduce, pues, a demostrar que la función $r$ es localmente constante. Para $m \in \Z$, se tiene $r(x) = m$ si y solo si $\uH^{m}(P')_{\ol{x}} \neq 0$; por tanto, en virtud de (SGAA~IX~2.4), $r$ es localmente constructible. Por consiguiente, para verificar que es localmente constante, podemos, por (EGA~IV~1.10.1), suponer que $X$ es local. Sea $x$ el punto cerrado de $X$, y sea $U$ el abierto complementario.

Vamos a demostrar que $r$ es constante. Razonando por inducción sobre la dimensión de $X$, podemos suponer ya que $r(y)$ es constante e igual a $r_0$ para $y \in U$, y se trata de probar que $r(x) = r_0$.

Supongamos que no es así, es decir, que $r(x) \neq r_0$, y demostremos que se llega a una contradicción. Por hipótesis, se tiene $P|U \xrightarrow{\sim} L[r_0]$, donde $L$ es un $A_U$-módulo invertible. Trasladando los grados si es necesario, podemos suponer que $r_0 = 0$. Por otra parte, se tiene $r(x) \neq 0$, digamos $r(x) = a > 0$. Denotemos por $i : U \to X$ y $j : x \to X$ las inclusiones canónicas.

Escribamos la sucesión exacta:
\[
(**) \quad 0 \longrightarrow i_{!}(P|U) \longrightarrow P \longrightarrow j_{*}j^{*}P \longrightarrow 0.
\]
Por hipótesis, los únicos objetos de cohomología no nulos de $P$ son $\uH^{0}(P)$ y $\uH^{-a}(P)$, y se tiene $\uH^{0}(P) \xrightarrow{\sim} i_{!}(L)$. De ello se sigue que $(**)$ se escinde naturalmente (gracias a la flecha $P \to \uH^{0}(P)$), de modo que se tiene
\[
P \xrightarrow{\sim} i_{!}(P|U) + j_{*}j^{*}P.
\]
Se deduce de ello
\begin{align*}
P \mathop{\otimes}\limits^{\mathrm{L}} P' &\xrightarrow{\sim} i_{!}(P|U) \mathop{\otimes}\limits^{\mathrm{L}} P' + j_{*}j^{*}P \mathop{\otimes}\limits^{\mathrm{L}} P' \\
&\xrightarrow{\sim} i_{!}(P|U \mathop{\otimes}\limits^{\mathrm{L}} P'|U) + j_{*}(j^{*}P \mathop{\otimes}\limits^{\mathrm{L}} j^{*}P') \\
&\xrightarrow{\sim} i_{!}(A_U) + j_{*}(A_x),
\end{align*}
lo cual es absurdo, puesto que $i_{!}(A_U) + j_{*}(A_x)$ manifiestamente no es isomorfo a $A_X$ (¡ya los $\uH^{0}$ no son isomorfos!).

Esto concluye la demostración del lema 2.2.
